\documentclass[UTF8]{ctexart}

\usepackage[a4paper,margin=2.5cm]{geometry}
\usepackage{amsmath,amssymb}
\usepackage{booktabs}
\usepackage{enumitem}
\usepackage{hyperref}

\title{LPA* 动态路径规划算法说明与方法论分析}
\author{运筹优化课程项目}
\date{}

\begin{document}
\maketitle

\section{问题定位：从通用搜索到动态运筹优化}

迷宫路径规划可以被抽象为图论中的最短路问题。给定图
$G=(V,E)$、起点 $s$、终点集合 $T$，静态无权情形只需最小化路径边数：
\[
    \min_{P:s\leadsto T} |P|-1 .
\]
这对应 BFS 的问题类：约束少、适用面广，但目标函数表达能力弱。A* 在 BFS
的基础上引入启发函数 $h(v)$，用 $f(v)=g(v)+h(v)$ 改变搜索顺序；它仍然求解同一类静态最短路问题，只是用 admissible heuristic
提高搜索效率。Weighted A* 和 Cost-aware A* 进一步把启发权重、转弯、风险区、狭窄区等因素加入模型，问题类变窄，但目标函数更接近真实机器人运行。

本项目中的 LPA* 并不只追求“跑出一条路径”，而是面向一类更具体的问题：

\begin{quote}
在地图结构基本稳定、局部障碍会随时间变化、机器人需要兼顾长度、转弯、风险暴露和狭窄区域通过成本的动态路径规划问题。
\end{quote}

这体现了一个运筹优化中的常见取舍：算法越通用，往往越难表达复杂目标；算法越贴合业务结构，适用问题类越小，但优化质量和解释力越强。LPA* 的价值正在于这种中间位置：它比静态 A* 更适合动态变化，又比完全重新规划更能复用历史计算。

\section{图模型增强：方向状态与多目标边权}

普通网格图以格子 $(x,y)$ 为节点。但如果目标函数包含转弯成本，仅用格子节点是不够的，因为“到达同一格子”的方向不同，会影响下一步转弯代价。因此程序将状态空间扩展为
\[
    z=(x,y,d_x,d_y),
\]
其中 $(d_x,d_y)$ 表示进入当前格子的方向。起点使用无方向状态 $(x_s,y_s,\varnothing,\varnothing)$。

对从状态 $z_i$ 移动到 $z_j$ 的边，定义代价：
\[
    c(z_i,z_j)
    =
    p_1
    + p_2 \cdot \tau(z_i,z_j)
    + p_3 \cdot I_{\mathrm{risk}}(v_j)
    + p_4 \cdot I_{\mathrm{narrow}}(v_j),
\]
其中：

\begin{itemize}[leftmargin=2em]
    \item $p_1$ 是单位步长成本；
    \item $\tau(z_i,z_j)$ 是转弯成本，直行取 $0$，直角转弯取 $1$，掉头取 $2$；
    \item $I_{\mathrm{risk}}(v_j)$ 表示是否进入风险格；
    \item $I_{\mathrm{narrow}}(v_j)$ 表示是否进入狭窄格。
\end{itemize}

最终报告的综合代价为：
\[
    J(P)
    =
    p_1 L(P)
    + p_2 R(P)
    + p_3 Q(P)
    + p_4 N(P)
    + p_5 K(P),
\]
其中 $L$ 为路径长度，$R$ 为转弯次数，$Q$ 为风险区通过次数，$N$ 为狭窄区通过次数，$K$ 为重规划次数。前四项进入图上的边权，最后一项用于评价动态修复过程的稳定性。

这种建模使 LPA* 从“动态最短路算法”变成“动态多成本最短路算法”。它牺牲了一部分通用性，例如状态空间从 $|V|$ 扩展到约 $5|V|$，但换来了对实际机器人运行质量的更强表达。

\section{LPA* 的核心思想：局部一致性与增量修复}

LPA* 维护两个值：

\[
    g(z): \text{当前已确认的最优估计}, \qquad
    rhs(z): \text{由后继节点诱导的一步前瞻估计}.
\]

在目标根搜索的实现中，终点区域状态满足
\[
    rhs(t)=0, \quad t\in T,
\]
其他状态初始为无穷。对非目标状态：
\[
    rhs(z)=\min_{z'\in Succ(z)} \left(c(z,z') + g(z')\right).
\]

如果 $g(z)=rhs(z)$，称节点 $z$ 局部一致；如果不相等，则说明它可能受到了地图变化或邻居值变化的影响，需要进入优先队列。队列键值为：
\[
    key(z)=
    \left(
    \min(g(z),rhs(z)) + h(s,z),
    \min(g(z),rhs(z))
    \right).
\]

这里 $h(s,z)$ 使用乘以 $p_1$ 的曼哈顿距离，作为 admissible 的低估启发。动态障碍发生时，程序只更新变化格及其邻接格对应的方向状态，然后通过队列传播影响。这比从头运行 A* 更符合“局部扰动、局部修复”的运筹思想。

\section{与已有算法的比较}

\begin{table}[h]
\centering
\begin{tabular}{llll}
\toprule
算法 & 适用问题类 & 优化目标 & 动态变化处理 \\
\midrule
BFS & 静态无权图 & 最短步数 & 不支持，需要重跑 \\
DFS & 静态可达性 & 找到一条路 & 不保证最优 \\
A* & 静态启发式最短路 & 最短步数，少扩展 & 不支持，需要重跑 \\
Cost-aware A* & 静态多成本路径 & 长度、转弯、风险、狭窄 & 不支持，需要重跑 \\
Weighted A* & 静态加权启发式路径 & 更快但可能牺牲最优性 & 不支持，需要重跑 \\
LPA* & 局部动态多成本路径 & 多成本最短路和增量修复 & 支持局部更新 \\
D* Lite & 在线动态路径修正 & 机器人运动中的增量修复 & 支持局部更新 \\
\bottomrule
\end{tabular}
\end{table}

从方法论上看，BFS 和 A* 解决的是“最短”问题；Cost-aware A* 解决的是“更合理”问题；LPA* 解决的是“环境变化后如何少算而仍然合理”的问题。它的思维深度不在于单次路径一定比所有算法短，而在于利用问题的时间结构：相邻时刻的最优路径通常高度相关，因此不应丢弃上一轮计算。

\section{程序实现对应关系}

程序保持统一接口：
\[
    \texttt{run\_lpa(map\_input, dynamic\_updates=None, **kwargs) -> dict}.
\]

本次实现保留原有返回字段，并额外增强以下信息：

\begin{itemize}[leftmargin=2em]
    \item \texttt{cost\_weights}: 记录 $p_1,\ldots,p_5$，使实验可复现；
    \item \texttt{optimization\_profile}: 说明状态空间、目标函数和适用问题类；
    \item \texttt{explored\_cells}: 记录搜索展开的格子，用于可视化；
    \item \texttt{updated\_nodes}: 统计动态事件实际影响到的方向状态数；
    \item \texttt{changed\_cells}: 保留原始变化格数量，用于区分环境扰动规模和算法修复规模；
    \item \texttt{visualization\_trace}: 记录初始规划、动态更新、修复和最终路径。
\end{itemize}

这些字段不破坏主程序接口，因为注册器会保留统一字段，并允许算法结果包含额外解释性字段。

\section{适用性、局限与进一步改进}

LPA* 的优势在于动态局部变化。当地图每一步都大范围改变时，增量修复的优势会下降，甚至可能不如直接重跑 A* 简洁。因此它适合“变化稀疏、路径目标稳定、代价结构明确”的场景。

本实现的主要局限有三点：

\begin{enumerate}[leftmargin=2em]
    \item 状态空间扩大到方向状态后，理论上会增加内存和队列操作成本；
    \item 动态事件当前只处理阻塞和释放，没有处理风险区、狭窄区权重随时间变化；
    \item 启发函数只使用距离下界，没有纳入可证明 admissible 的转弯下界。
\end{enumerate}

进一步可以研究三类扩展：第一，引入风险概率，把 $I_{\mathrm{risk}}$ 改为期望损失；第二，引入可调权重的参数敏感性分析，观察路径对 $p_1,\ldots,p_5$ 的响应；第三，把 LPA* 与 D* Lite 的实验结果放在同一动态事件序列下比较，分析“固定起点增量规划”和“在线移动修正”的差别。

\section{结论}

本项目中的 LPA* 体现了从算法到运筹模型的转换：它不是简单地把 A* 改成动态版本，而是将状态、边权、动态扰动和评价指标统一到图论优化框架中。其核心取舍是明确的：用较窄的问题类换取更强的目标表达和更高的动态复用能力。这样的设计更接近实际路径规划中的决策需求，也能与 BFS、A*、Weighted A* 形成清晰的层次比较。

\end{document}
